% Part: computability
% Chapter: recursive-functions
% Section: trees

\documentclass[../../../include/open-logic-section]{subfiles}

\begin{document}

\olfileid[tr]{cmp}{rec}{tre}
\olsection{Ağaçlar}

Dizileri sayılarla, diziler üzerindeki özellik ve işlemleri de kodlar
üzerindeki ilkel özyinelemeli bağıntı ve fonksiyonlarla temsil edebildiğimiz
gibi, bazen ağaçları da doğal sayılar olarak temsil etmek yararlıdır. Bunun
için dizileri ve kodlarını kullanacağız. Bir ağaç, tek bir düğümden (belki
etiketli) ya da belirli sayıda alt ağaca bağlı bir düğümden (yine belki
etiketli) oluşabilir. Bu düğüme ağacın \emph{kökü}, bağlı olduğu alt ağaçlara
ise ağacın \emph{dolaysız alt ağaçları} denir.

Ağaçları özyinelemeli olarak $\tuple{k,d_1,\dots,d_k}$ dizisiyle kodlarız;
burada $k$ dolaysız alt ağaçların sayısı, $d_1,\dots,d_k$ ise bu alt ağaçların
kodlarıdır. Düğümlerin etiketleri varsa dolaysız alt ağaçlardan sonra
eklenebilir. Böylece yalnızca $l$ etiketli tek düğümden oluşan ağaç
$\tuple{0,l}$ ile; $l_1$ etiketli bir kökün $l_2$ ve $l_3$ etiketli iki tek
düğüme bağlandığı ağaç ise
$\tuple{2,\tuple{0,l_2},\tuple{0,l_3},l_1}$ ile kodlanır.

\begin{prop}
  \ollabel{prop:subtreeseq}
  Kodu $t$ olan ağacın bütün alt ağaçlarının kodlarını eleman olarak içeren
  dizinin kodunu döndüren $\fn{SubtreeSeq}(t)$ fonksiyonu ilkel
  özyinelemelidir.
\end{prop}

\begin{proof}
  Önce $\fn{ISubtrees}(t)=\fn{subseq}(t,1,(t)_0)$ fonksiyonunun ilkel
  özyinelemeli olduğuna ve $t$ ağacının dolaysız alt ağaçlarının kodlarını
  döndürdüğüne dikkat edin. Şimdi kökten $n$ düğüm uzakta bulunan bütün alt
  ağaçların dizisini hesaplayan yardımcı bir $\fn{hSubtreeSeq}(t,n)$
  fonksiyonu tanımlayabiliriz. Kökten $0$ düğüm uzakta bulunan, başka bir
  deyişle $t$'nin kökünde başlayan alt ağaçlar dizisi yalnızca $t$'den
  oluşur. $t$'nin $(n+1)$. düzeydeki bütün alt ağaçlarının dizisini elde
  etmek için $n$. düzeydeki alt ağaçların dizisini, bunların bütün dolaysız
  alt ağaçlarından oluşan diziyle art arda birleştiririz. Bunların tümünün
  listesini elde etmek için şu gözlemi kullanın: $f(x)$ dizi kodları döndüren
  ilkel özyinelemeli bir fonksiyonsa
  $g_f(s,k)=f((s)_0)\concat\dots\concat f((s)_k)$ de ilkel
  özyinelemelidir:
  \begin{align*}
    g(s,0) & = f((s)_0)\\
    g(s,k+1) & = g(s,k)\concat f((s)_{k+1}).
  \end{align*}
  Örneğin $s$ bir ağaç dizisiyse
  $h(s)=g_{\fn{ISubtrees}}(s,\len{s})$, $s$ elemanlarının dolaysız alt
  ağaçları dizisini verir. Bunu kullanarak $\fn{hSubtreeSeq}$ fonksiyonunu
  şöyle tanımlarız:
  \begin{align*}
    \fn{hSubtreeSeq}(t,0) & = \tuple{t}\\
    \fn{hSubtreeSeq}(t,n+1) & =
      \fn{hSubtreeSeq}(t,n)\concat h(\fn{hSubtreeSeq}(t,n)).
  \end{align*}
  Kodu $t$ olan bir ağaçtaki en büyük alt ağaç düzeyi, yani kökle bir yaprak
  düğüm arasındaki en büyük uzaklık $t$ koduyla sınırlıdır. Dolayısıyla kodu
  $t$ olan ağacın bütün alt ağaçlarının kodları dizisi
  $\fn{hSubtreeSeq}(t,t)$ ile verilir.
\end{proof}

\begin{prob}
  \olref[cmp][rec][tre]{prop:subtreeseq} ispatındaki
  $\fn{hSubtreeSeq}$ tanımı genel olarak tekrarlar içerir. Bir alt ağacın
  kodunun sonuç listesindeki yalnızca bir kez geçmesini güvenceye alan
  alternatif bir tanım verin.
\end{prob}

\end{document}

